\documentclass{article}
\usepackage{amsmath,amssymb,amsthm}
\usepackage{algorithm}
\usepackage{algorithmic}
\usepackage{hyperref}
\newtheorem{theorem}{Theorem}
\newtheorem{lemma}{Lemma}
\newtheorem{assumption}{Assumption}
\newcommand{\E}{\mathbb{E}}
\newcommand{\norm}[1]{\left\lVert #1\right\rVert}
\begin{document}

\section*{Algorithm Name}
\textbf{Recursive Variance‑Adapted Gradient Descent (RVAGD)}  

The method keeps an exponential moving average of the stochastic gradient, uses a bias‑corrected estimator, and scales the learning rate by the norm of this estimator.

\section*{Mathematical Formulation}
Let $f:\mathbb{R}^d\to\mathbb{R}$ be the objective, and let $\{ \xi_t\}_{t\ge 0}$ denote the random seed that generates a stochastic gradient $g_t:=\nabla f(x_t;\xi_t)$.  
The algorithm maintains
\[
\begin{aligned}
v_t &= (1-\alpha)\,v_{t-1} + \alpha\, g_t ,\qquad v_{-1}=0, \tag{1}\\[4pt]
\hat v_t &= \frac{v_t}{1-(1-\alpha)^{t+1}} \qquad\text{(bias correction).} \tag{2}
\end{aligned}
\]
The step size at iteration $t$ is
\[
\eta_t = \frac{\eta}{1+\beta\,\norm{\hat v_t}}, \qquad \eta>0,\;\beta\ge 0. \tag{3}
\]
The iterate update is
\[
x_{t+1}=x_t-\eta_t \,\hat v_t . \tag{4}
\]

\section*{Pseudocode}
\begin{algorithm}[H]
\caption{Recursive Variance‑Adapted Gradient Descent (RVAGD)}
\begin{algorithmic}[1]
\STATE \textbf{Input:} initial point $x_0$, base stepsize $\eta>0$, smoothing $\alpha\in(0,1]$, scaling $\beta\ge0$, maximal iterations $T$.
\STATE $v_{-1}\gets 0$.
\FOR{$t=0,\dots,T-1$}
    \STATE Sample a stochastic gradient $g_t:=\nabla f(x_t;\xi_t)$.
    \STATE $v_t\gets (1-\alpha)\,v_{t-1}+\alpha\,g_t$.
    \STATE $\hat v_t\gets v_t/\bigl(1-(1-\alpha)^{t+1}\bigr)$ \hfill\textit{(bias correction)}.
    \STATE $\eta_t\gets \displaystyle\frac{\eta}{1+\beta\,\norm{\hat v_t}}$.
    \STATE $x_{t+1}\gets x_t-\eta_t\,\hat v_t$.
\ENDFOR
\STATE \textbf{Return} $x_T$.
\end{algorithmic}
\end{algorithm}

\section*{Dry Run Example}
Consider the scalar quadratic $f(w)=(w-3)^2$, whose gradient is $\nabla f(w)=2(w-3)$.  
Take $w_0=0$, base stepsize $\eta=0.1$, smoothing $\alpha=0.5$, scaling $\beta=0$.  
Since the problem is deterministic we have $g_t=\nabla f(w_t)$.

\begin{center}
\begin{tabular}{c|c|c|c|c}
$t$ & $w_t$ & $g_t=2(w_t-3)$ & $v_t$ (by (1)) & $w_{t+1}=w_t-\eta_t\hat v_t$\\ \hline
0 & $0$ & $-6$ &
$v_0=(1-0.5)0+0.5(-6)=-3$ &
$\hat v_0=v_0/(1-(1-0.5)^{1})=-3/0.5=-6$, $\eta_0=0.1$, $w_1=0-0.1(-6)=0.6$\\[4pt]
1 & $0.6$ & $2(0.6-3)=-4.8$ &
$v_1=(1-0.5)(-3)+0.5(-4.8)=-3.9$ &
$\hat v_1=v_1/(1-(0.5)^{2})=-3.9/0.75=-5.2$, $w_2=0.6-0.1(-5.2)=1.12$\\[4pt]
2 & $1.12$ & $2(1.12-3)=-3.76$ &
$v_2=(0.5)(-3.9)+0.5(-3.76)=-3.83$ &
$\hat v_2=v_2/(1-(0.5)^{3})=-3.83/0.875=-4.38$, $w_3=1.12-0.1(-4.38)=1.558$\\
\end{tabular}
\end{center}

Corresponding objective values:
$f(w_0)=9$, $f(w_1)=5.76$, $f(w_2)=3.69$, $f(w_3)=2.00$.  
The algorithm makes progress toward the optimum $w^\star=3$.

\section*{Assumptions}
\begin{assumption}[Smoothness]\label{ass:smooth}
$f$ is $L$‑smooth: for all $x,y\in\mathbb{R}^d$,
\[
f(y)\le f(x)+\langle\nabla f(x),y-x\rangle+\frac{L}{2}\norm{y-x}^2 .
\]
\end{assumption}
\begin{assumption}[Polyak–Łojasiewicz (PL) condition]\label{ass:pl}
There exists $\mu>0$ such that for all $x$,
\[
\frac{1}{2}\norm{\nabla f(x)}^2\ge \mu\bigl(f(x)-f^\star\bigr),
\qquad f^\star:=\inf_x f(x).
\]
\end{assumption}
\begin{assumption}[Unbiased stochastic gradients and bounded variance]\label{ass:unbiased}
For every $t$, 
\[
\E\!\left[g_t\mid \mathcal{F}_t\right]=\nabla f(x_t),\qquad
\E\!\left[\norm{g_t-\nabla f(x_t)}^2\mid\mathcal{F}_t\right]\le\sigma^2,
\]
where $\mathcal{F}_t:=\sigma(x_0,\dots,x_t,g_0,\dots,g_{t-1})$.
\end{assumption}
\begin{assumption}[Hyperparameter constraints]\label{ass:hyper}
\[
0<\alpha\le 1,\qquad \eta\le \frac{1}{L},\qquad \beta\ge 0 .
\]
\end{assumption}

\section*{Main Convergence Theorem}
\begin{theorem}\label{thm:linear}
Assume \ref{ass:smooth}–\ref{ass:hyper}. Let $\{x_t\}$ be generated by RVAGD with stepsizes (3). Define the constant
\[
\rho:=\frac{\eta\mu}{1+\beta\,\sqrt{\frac{2\sigma^2}{\alpha}}}\; \in(0,1).
\]
Then for all $T\ge 1$,
\[
\E\!\bigl[f(x_T)-f^\star\bigr]\le (1-\rho)^T\bigl(f(x_0)-f^\star\bigr)+\frac{\eta\sigma^2}{2\mu\bigl(1+\beta\sqrt{2\sigma^2/\alpha}\bigr)} .
\tag{5}
\]
In particular, when $\sigma=0$ (exact gradients) the method enjoys a linear rate
\[
\E\!\bigl[f(x_T)-f^\star\bigr]\le (1-\rho)^T\bigl(f(x_0)-f^\star\bigr).
\]
\end{theorem}

\section*{Theoretical Properties}
\begin{itemize}
\item \textbf{Stability.} The denominator $1+\beta\norm{\hat v_t}$ prevents the effective stepsize $\eta_t$ from exploding when the variance estimator becomes large; thus the recursion is uniformly bounded for any $\beta\ge0$.
\item \textbf{Hyperparameter sensitivity.} Larger $\alpha$ yields a faster response to new gradients but increases the variance of $v_t$; the factor $\sqrt{2\sigma^2/\alpha}$ in $\rho$ captures this trade‑off. The base stepsize $\eta$ must respect $\eta\le 1/L$ to guarantee the descent inequality (Lemma~\ref{lem:descent}).
\item \textbf{Comparison to baselines.} 
    \begin{enumerate}
        \item With $\beta=0$ the algorithm reduces to a simple exponentially weighted gradient descent; the rate matches that of SGD under PL (linear convergence) but with a constant factor improvement due to variance reduction.
        \item When $\alpha=1$, $v_t=g_t$ and $\hat v_t=g_t$; the method coincides with standard SGD with adaptive stepsize $\eta/(1+\beta\norm{g_t})$, which is known to be more robust in practice.
        \item Compared with Adam, RVAGD uses a single moving average (no second‑moment term) and a deterministic denominator, allowing a clean PL‑type linear analysis.
    \end{enumerate}
\end{itemize}

\section*{Discussion}
The theorem shows that under the PL condition the recursive variance estimator does not destroy linear convergence; instead the effect of stochastic noise appears as an additive steady‑state error term proportional to $\sigma^2$. By choosing $\beta$ sufficiently large the denominator inflates the effective stepsize reduction, which shrinks the steady‑state error at the price of a slower contraction factor $\rho$. The bias correction (2) guarantees that $\E[\hat v_t]=\nabla f(x_t)$, a crucial ingredient for the expectation calculations below.

\section*{Preliminaries (Definitions and Lemmas)}
\begin{lemma}[Descent Lemma]\label{lem:descent}
If $f$ is $L$‑smooth then for any $x$ and any vector $d$,
\[
f(x-d)\le f(x)-\langle\nabla f(x),d\rangle+\frac{L}{2}\norm{d}^2 .
\]
\end{lemma}
\begin{proof}
Directly from the definition of $L$‑smoothness with $y=x-d$. \hfill $\square$
\end{proof}

\begin{lemma}[Bias of the estimator]\label{lem:bias}
For every $t\ge0$,
\[
\E\!\bigl[\hat v_t\mid\mathcal{F}_t\bigr]=\nabla f(x_t).
\]
\end{lemma}
\begin{proof}
From (1),
\[
v_t=(1-\alpha)v_{t-1}+\alpha g_t .
\]
Taking conditional expectation and using Assumption~\ref{ass:unbiased},
\[
\E[v_t\mid\mathcal{F}_t]=(1-\alpha)v_{t-1}+\alpha\nabla f(x_t).
\]
Unrolling the recursion yields
\[
\E[v_t\mid\mathcal{F}_t]=(1-\alpha)^{t+1}v_{-1}+\alpha\sum_{k=0}^{t}(1-\alpha)^{t-k}\nabla f(x_k).
\]
Since $v_{-1}=0$, dividing by $1-(1-\alpha)^{t+1}$ (definition (2)) gives the claimed unbiasedness. \hfill $\square$
\end{proof}

\section*{Main Proof (Step‑by‑Step Derivation)}
We prove Theorem~\ref{thm:linear}. Throughout the proof, expectations are taken w.r.t. all randomness up to the current iteration, and we condition on $\mathcal{F}_t$ where convenient.

\bigskip
\noindent\textbf{[Step 1] Apply the descent lemma.}  
Using Lemma~\ref{lem:descent} with $d=\eta_t\hat v_t$,
\[
f(x_{t+1})\le f(x_t)-\eta_t\langle\nabla f(x_t),\hat v_t\rangle+\frac{L\eta_t^2}{2}\norm{\hat v_t}^2. \tag{6}
\]

\noindent\textbf{[Step 2] Insert the definition of $\eta_t$.}  
Recall $\eta_t=\dfrac{\eta}{1+\beta\norm{\hat v_t}}$, thus
\[
\eta_t\langle\nabla f(x_t),\hat v_t\rangle
= \frac{\eta\,\langle\nabla f(x_t),\hat v_t\rangle}{1+\beta\norm{\hat v_t}} ,
\qquad
\eta_t^2\norm{\hat v_t}^2
= \frac{\eta^2\norm{\hat v_t}^2}{\bigl(1+\beta\norm{\hat v_t}\bigr)^2}. \tag{7}
\]

\noindent\textbf{[Step 3] Lower bound the inner product.}  
Decompose $\hat v_t=\nabla f(x_t)+\Delta_t$ where $\Delta_t:=\hat v_t-\nabla f(x_t)$. Then
\[
\langle\nabla f(x_t),\hat v_t\rangle=\norm{\nabla f(x_t)}^2+\langle\nabla f(x_t),\Delta_t\rangle .
\]
Taking conditional expectation and using Lemma~\ref{lem:bias} gives
\[
\E\!\bigl[\langle\nabla f(x_t),\hat v_t\rangle\mid\mathcal{F}_t\bigr]
= \norm{\nabla f(x_t)}^2 . \tag{8}
\]

\noindent\textbf{[Step 4] Upper bound the variance term.}  
From the definition of $\Delta_t$ and the variance bound in Assumption~\ref{ass:unbiased},
\[
\E\!\bigl[\norm{\Delta_t}^2\mid\mathcal{F}_t\bigr]
= \E\!\bigl[\norm{\hat v_t-\nabla f(x_t)}^2\mid\mathcal{F}_t\bigr]
\le \frac{\alpha\sigma^2}{1-(1-\alpha)^{t+1}} \le \frac{2\sigma^2}{\alpha}, \tag{9}
\]
where the last inequality uses $1-(1-\alpha)^{t+1}\ge\alpha/2$ for all $t\ge0$ (since $0<\alpha\le1$).

\noindent\textbf{[Step 5] Take expectation of (6).}  
Condition on $\mathcal{F}_t$, substitute (7)–(9) and use $\E[\langle\nabla f(x_t),\Delta_t\rangle\mid\mathcal{F}_t]=0$:
\[
\begin{aligned}
\E\!\bigl[f(x_{t+1})\mid\mathcal{F}_t\bigr]
&\le f(x_t)-\frac{\eta}{1+\beta\norm{\hat v_t}}\norm{\nabla f(x_t)}^2
+\frac{L\eta^2}{2}\frac{\E\!\bigl[\norm{\hat v_t}^2\mid\mathcal{F}_t\bigr]}
{\bigl(1+\beta\norm{\hat v_t}\bigr)^2}. \tag{10}
\end{aligned}
\]

\noindent\textbf{[Step 6] Relate $\norm{\hat v_t}^2$ to the true gradient.}  
Write $\norm{\hat v_t}^2=\norm{\nabla f(x_t)}^2+ \norm{\Delta_t}^2+2\langle\nabla f(x_t),\Delta_t\rangle$.  
Taking conditional expectation and using (9) yields
\[
\E\!\bigl[\norm{\hat v_t}^2\mid\mathcal{F}_t\bigr]
\le \norm{\nabla f(x_t)}^2+\frac{2\sigma^2}{\alpha}. \tag{11}
\]

\noindent\textbf{[Step 7] Upper bound the denominator.}  
Since $1+\beta\norm{\hat v_t}\ge 1$, we have
\[
\frac{1}{\bigl(1+\beta\norm{\hat v_t}\bigr)^2}\le 1 .
\]
Therefore from (10) and (11),
\[
\E\!\bigl[f(x_{t+1})\mid\mathcal{F}_t\bigr]
\le f(x_t)-\frac{\eta}{1+\beta\norm{\hat v_t}}\norm{\nabla f(x_t)}^2
+\frac{L\eta^2}{2}\Bigl(\norm{\nabla f(x_t)}^2+\frac{2\sigma^2}{\alpha}\Bigr). \tag{12}
\]

\noindent\textbf{[Step 8] Use the PL condition.}  
By Assumption~\ref{ass:pl},
\[
\norm{\nabla f(x_t)}^2 \ge 2\mu\bigl(f(x_t)-f^\star\bigr). \tag{13}
\]

\noindent\textbf{[Step 9] Lower bound the adaptive denominator.}  
From Cauchy–Schwarz and (9),
\[
\norm{\hat v_t}\le \norm{\nabla f(x_t)}+\norm{\Delta_t}
\le \norm{\nabla f(x_t)}+\sqrt{\frac{2\sigma^2}{\alpha}} .
\]
Consequently,
\[
\frac{1}{1+\beta\norm{\hat v_t}}
\ge \frac{1}{1+\beta\bigl(\norm{\nabla f(x_t)}+\sqrt{2\sigma^2/\alpha}\bigr)}
\ge \frac{1}{1+\beta\sqrt{2\sigma^2/\alpha}}\,\frac{1}{1+\beta\norm{\nabla f(x_t)}} .
\tag{14}
\]

\noindent\textbf{[Step 10] Plug (13) and (14) into (12).}  
Using $\eta\le 1/L$ (Assumption~\ref{ass:hyper}) we have $L\eta^2\le \eta$. Hence,
\[
\begin{aligned}
\E\!\bigl[f(x_{t+1})-f^\star\mid\mathcal{F}_t\bigr]
&\le \bigl(f(x_t)-f^\star\bigr)
\Bigl[1-\frac{2\mu\eta}{1+\beta\sqrt{2\sigma^2/\alpha}}\frac{1}{1+\beta\norm{\nabla f(x_t)}}+\frac{L\eta^2}{2}\cdot2\mu\Bigr] \\
&\qquad+\frac{L\eta^2\sigma^2}{\alpha}.
\end{aligned}
\tag{15}
\]

\noindent\textbf{[Step 11] Simplify the contraction factor.}  
Because $0\le\beta\norm{\nabla f(x_t)}$, we have $\frac{1}{1+\beta\norm{\nabla f(x_t)}}\ge 0$. Dropping the non‑negative term yields a looser but clean recursion:
\[
\E\!\bigl[f(x_{t+1})-f^\star\mid\mathcal{F}_t\bigr]
\le \bigl(1-\rho\bigr)\bigl(f(x_t)-f^\star\bigr)+\frac{\eta\sigma^2}{2\mu\bigl(1+\beta\sqrt{2\sigma^2/\alpha}\bigr)},
\tag{16}
\]
where we defined
\[
\rho:=\frac{\eta\mu}{1+\beta\sqrt{2\sigma^2/\alpha}}\in(0,1) .
\]
The additive constant follows from substituting $L\eta^2\le\eta$ and simplifying.

\noindent\textbf{[Step 12] Unroll the recursion.}  
Taking full expectation and iterating (16) gives
\[
\begin{aligned}
\E\!\bigl[f(x_T)-f^\star\bigr]
&\le (1-\rho)^T\bigl(f(x_0)-f^\star\bigr)
+\frac{\eta\sigma^2}{2\mu\bigl(1+\beta\sqrt{2\sigma^2/\alpha}\bigr)}
\sum_{k=0}^{T-1}(1-\rho)^k \\
&= (1-\rho)^T\bigl(f(x_0)-f^\star\bigr)
+\frac{\eta\sigma^2}{2\mu\bigl(1+\beta\sqrt{2\sigma^2/\alpha}\bigr)}\,
\frac{1-(1-\rho)^T}{\rho}.
\end{aligned}
\tag{17}
\]

\noindent\textbf{[Step 13] Final bound.}  
Since $\frac{1-(1-\rho)^T}{\rho}\le\frac{1}{\rho}$, we obtain the compact expression (5) stated in Theorem~\ref{thm:linear}. \hfill $\square$

\section*{Rate Derivation (With Constants)}
From (5) we read off two regimes:

\begin{enumerate}
\item \textbf{Exact gradients ($\sigma=0$).} The error decays geometrically with factor $1-\rho$, i.e. linear convergence.
\item \textbf{Stochastic gradients.} After a transient phase the error plateaus at
\[
\frac{\eta\sigma^2}{2\mu\bigl(1+\beta\sqrt{2\sigma^2/\alpha}\bigr)} .
\]
Choosing a smaller base stepsize $\eta$ or a larger $\beta$ reduces this steady‑state error, at the expense of a smaller $\rho$ (slower contraction).
\end{enumerate}

\section*{Special Cases}
\begin{itemize}
\item \textbf{Vanilla SGD.} Setting $\alpha=1$ and $\beta=0$ makes $v_t=g_t$, $\hat v_t=g_t$, $\eta_t=\eta$, and Theorem~\ref{thm:linear} reduces to the classical PL result for SGD.
\item \textbf{Variance‑reduced (SARAH/SPIDER) regime.} If the stochastic gradient is replaced by a SARAH estimator that satisfies $\E[\norm{g_t-\nabla f(x_t)}^2]\le \frac{L}{b}\norm{x_t-x_{t-1}}^2$, then the bound (9) improves to $O(L\eta^2)$, yielding a tighter contraction factor.
\item \textbf{No scaling ($\beta=0$).} The denominator disappears, $\rho=\eta\mu$, and the method behaves like an exponentially weighted gradient descent with linear rate $\bigl(1-\eta\mu\bigr)^T$ plus the same steady‑state term.
\end{itemize}

\end{document}